\begin{textsample}{POS Dim 2 – 2019\_c4 – Score 18.00 – 2019\_c4/enade\_2019\_tec\_agro\_85543dcc5c.txt}  \label{ex:f2_pos_042}
CAPÍTULO 4 PERCEPÇÃO DA PROVA Nas análises feitas neste capítulo, constam as percepções dos concluintes da Área de Tecnologia em Agronegócios sobre a prova aplicada no Enade/2019.
Essas percepções foram mensuradas por meio de nove questões que avaliaram desde o grau de dificuldade da prova até o tempo gasto para concluí -la.
As percepções sobre a prova foram relacionadas ao desempenho dos estudantes, à Grande Região de funcionamento do curso, à Categoria Administrativa e ao tipo de Organização Acadêmica da IES.
O \textbf{questionário} de percepção da prova encontra -se ao final do Anexo VII, que traz a \textbf{reprodução} do exame.
O desempenho dos estudantes foi classificado em quatro quartos.
Para tanto, esse desempenho foi \textbf{ordenado} de forma ascendente.
O percentil 25, P25, também conhecido como primeiro quartil, é a nota de desempenho que deixa um quarto ( 25 \% ) dos valores observados abaixo e três quartos acima.
A Figura 1 apresenta uma ilustração desse conceito.
O quarto inferior de desempenho é composto pelas notas abaixo do primeiro quartil.
Já o percentil 75, P75, também conhecido como terceiro quartil, é o valor para o qual há três quartos ( 75 \% ) dos dados abaixo e um quarto acima dele.
O quarto superior de desempenho é composto pelas notas iguais ou acima do terceiro quartil.
O percentil 50, P50, também conhecido como mediana, é o valor que divide as notas em dois conjuntos de igual tamanho.
O segundo quarto inclui valores entre o primeiro quartil ( P25 ) e a mediana.
O terceiro quarto contém os valores entre a mediana ( P50 ) e o terceiro quartil ( P75 ).
Vale ressaltar que percentis, quartis e medianas são pontos que, não obrigatoriamente, pertencem ao conjunto original de dados, ao passo que os quartos são \textbf{subconjuntos} dos dados originais.
A seguir, serão apresentados gráficos com resultados selecionados, relativos às nove questões avaliadas por grupos de estudantes.
Nas barras dos gráficos, apresenta -se o percentual de alunos que assinalaram uma das opções ou a soma das porcentagens daqueles que assinalaram duas ( ou três ) delas.
Por exemplo, para as questões 1 e 2, nos gráficos, é apresentada a porcentagem total de participantes que assinalaram as opções ( D ) Difícil e ( E ) Muito difícil.
Em cada barra, foram assinalados, também, os extremos do intervalo de confiança de 95 \% como linhas verticais unidas por uma linha horizontal na forma da letra H maiúscula.
O \textbf{estimador} de um \textbf{parâmetro} com um certo nível de confiança ( e.g. 95,0 \% ) deve conter o \textbf{parâmetro} no intervalo de confiança em 95 \% das vezes.
Na comparação entre os \textbf{estimadores} dos \textbf{parâmetros} de duas classes de uma dada categoria ( e.g.
Norte e Nordeste nas Grandes Regiões, de primeiro e último quarto dentro de desempenho, de IES Públicas e Privadas, ou de Universidades e Faculdades ) associados aos seus respectivos intervalos de confiança, diz -se que não há diferença estatisticamente significativa entre os \textbf{parâmetros} das duas categorias quando há \textbf{interseção} entre os intervalos de confiança e que há diferença, se os intervalos de confiança são \textbf{disjuntos} ( para maiores informações \textbf{vide} Glossário ).
Nas Tabelas no Anexo II, são apresentados os valores absolutos e a distribuição percentual18 das alternativas válidas das nove questões, segundo o mesmo \textbf{recorte} de desempenho dos alunos e Grande Região de funcionamento do curso, de Categoria Administrativa e do tipo de Organização Acadêmica da IES. 4.1 GRAU DE DIFICULDADE DA PROVA 4.1.1 Componente de Formação Geral Ao avaliarem Qual o grau de dificuldade desta prova na parte de Formação Geral? ( Questão 1 ), 40,0 \% do grupo de inscritos e presentes optaram pelas alternativas Difícil ou Muito difícil.
Entretanto, para mais da metade dos estudantes ( 55,8 \% ), o Componente de Formação Geral da prova foi considerado com grau de dificuldade Médio ( Gráfico 4.1, Gráfico 4.2 e, no Anexo II, a Tabela II.1 ).
O percentual de estudantes que consideraram a prova como Difícil ou Muito difícil foi maior na região Nordeste, onde a proporção foi de 44,0 \%, enquanto a de menor incidência foi a Centro-Oeste, com 38,6 \%.
No Gráfico 4.1, é possível observar que não existem diferenças estatisticamente significativas entre as regiões.
Nas Grandes Regiões, a proporção de presentes à prova que consideraram o Componente de Formação Geral como sendo de grau de dificuldade Médio esteve entre 50,0 \%, na região Nordeste e 58,0 \%, na região Centro-Oeste. 18 Cumpre lembrar uma das convenções para tabelas \textbf{numéricas} ( pág. iii ) sobre a possibilidade de a soma das partes não resultar em 100 \% por questões de \textbf{arredondamento}. de dificuldade Médio esteve entre 50,0 \%, na região Nordeste e 58,0 \%, na região Centro-Oeste.
O percentual de alunos que consideraram o componente de Formação Geral da prova Difícil ou Muito difícil foi crescente em função dos quartos de desempenho : 37,9 \%, no primeiro quarto e 45,1 \%, no quarto, grupo de melhor desempenho na prova.
Nos quartos de desempenho intermediários, a proporção de alunos que consideraram a prova Difícil ou Muito difícil foi de 38,8 \%, no segundo quarto e de 38,2 \%, no terceiro quarto.
As diferenças entre todos os quartos de desempenho não são estatisticamente significativas.
Para todos os quartos de desempenho, a alternativa modal para essa pergunta foi a Médio, com 50,0 \% e 58,2 \% dos respondentes nos quartos extremos, quarto e terceiro, respectivamente.
O percentual de estudantes que consideraram o componente de Formação Geral da prova como Difícil ou Muito difícil foi maior para os estudantes de CEFET/IF, para os quais a proporção foi de 61,1 \%, enquanto os de Universidades foram a Organização Acadêmica com a menor incidência, com 29,9 \%, ainda que sem diferença estatisticamente significativa de Universidades, com 51,8 \%.
No Gráfico 4.3, é possível observar que não existem diferenças estatisticamente significativas entre as Organizações Acadêmicas.
Ainda sobre essas organizações é possível observar que a proporção de presentes à prova que consideraram o Componente de Formação Geral como sendo de grau de dificuldade Médio esteve entre 51,8 \%, nas Universidades e 61,1 \%, nos CEFET/IF ( ver também a Tabela II.2, no Anexo II ).
O percentual de estudantes que consideraram o componente de Formação Geral da prova como Difícil ou Muito difícil foi maior para os estudantes de cursos de IES Privadas ( 45,5 \% ), uma diferença estatisticamente significativa da proporção em IES Públicas ( 36,5 \% ).
Já a proporção de presentes à prova que consideraram o Componente de Formação Geral como sendo de grau de dificuldade Médio foi 59,2 \% para os de IES Públicas e 50,4 \% para os de IES Privadas ( ver Gráfico 4.4 ). 4.1.2 Componente de Conhecimento Específico Ao responderem à Questão 2 – Qual o grau de dificuldade desta prova na parte de Componente Específico? – 34,4 \% do grupo de estudantes classificaram -na como Difícil ou Muito difícil.
Além disso, o Componente modal de Conhecimento Específico da prova foi considerado com grau de dificuldade Médio por 59,1 \% dos alunos ( Gráfico 4.5, Gráfico 4.6, e, no Anexo II, a Tabela II.3 ).
A análise das respostas dos estudantes quanto ao grau de dificuldade do Componente de Conhecimento Específico da prova, agregado por Grande Região, mostra que não existem diferenças estatisticamente significativas entre todas as proporções de alunos que a avaliaram como Difícil ou Muito difícil.
O percentual de alunos que classificaram o grau de dificuldade como Médio, no Componente de Conhecimento Específico, variou de 54,5 \% a 61,2 \%, para as regiões Centro-Oeste e Nordeste, respectivamente.
Considerando -se a avaliação da dificuldade das questões do Componente de Conhecimento Específico da prova, de acordo com o desempenho dos estudantes, observa -se que não há diferenças estatisticamente significativas dos resultados entre todos os quartos de desempenho.
A proporção dos que classificaram a parte específica como Difícil ou Muito difícil variou 33,7 \% ( terceiro quarto ) a 36,0 \% ( primeiro quarto ).
As demais proporções dessa resposta foram de 34,2 \% e 33,8 \%, nos segundo e quarto quartos, respectivamente.
Já a proporção dos que responderam que o grau de dificuldade das questões do Componente de Conhecimento Específico da prova foi Médio variou de 56,9 \%, no último quarto a 60,1 \%, no segundo.
Considerando -se a avaliação da dificuldade das questões do Componente de Conhecimento Específico da prova, o percentual de \textbf{estudantes} que as consideraram como Difícil ou Muito difícil foi maior para os estudantes de Universidades e Centros Universitários, para os quais a proporção, em ambos, foi de 35,6 \%, enquanto as Faculdades foram a Organização Acadêmica com a menor incidência, com 33,4 \%.
No Gráfico 4.7, é possível observar que não existem diferenças estatisticamente significativas entre todas as organizações.
Nas Organizações Acadêmicas, a proporção de presentes à prova que consideraram o Componente de Formação Geral como sendo de grau de dificuldade Médio esteve entre 56,1 \%, nas Universidades e 60,8 \% nos CEFET/IF ( ver também Tabela II.4, no Anexo II ).
O percentual de estudantes que consideraram o Componente Específico da prova como Difícil ou Muito difícil foi maior para os estudantes de cursos de IES Privadas ( 38,3 \% ), e não há diferença estatisticamente significativa da proporção em IES Públicas ( 31,9 \% ).
Já a proporção de presentes à prova que consideraram este componente como sendo de grau de dificuldade Médio foi 61,4 \% para os de IES Públicas e 55,6 \% para os de IES Privadas ( ver Gráfico 4.8 e Tabela II.4, no Anexo II ). 4.2 EXTENSÃO DA PROVA EM RELAÇÃO AO TEMPO TOTAL \textbf{Indagados} quanto à extensão da prova, em relação ao tempo total oferecido para a sua resolução ( Questão 3 ), os estudantes apontaram, com maior incidência, a alternativa que considerava a extensão Adequada, para todas as agregações consideradas ( Gráfico 4.9, Gráfico 4.10, e, no Anexo II, a Tabela II.5 ).
O percentual de alunos que responderam ser a extensão da prova adequada foi de 69,0 \%.
Já 26,7 \% dos inscritos presentes consideraram que a prova foi Longa ou Muito longa, e 4,2 \% a avaliaram como Curta ou Muito curta.
Dentre as Grandes Regiões, a proporção daqueles que avaliaram a prova como Longa ou Muito longa, em relação ao tempo total destinado à sua resolução, variou de 20,0 \%, na região Nordeste até 29,5 \%, na região Centro-Oeste.
Para essas categorias, não existem diferenças estatisticamente significativas entre as regiões.
Considerando -se o desempenho dos alunos, nota -se ainda que 65,1 \% consideraram a extensão da prova Adequada, no quarto de desempenho inferior e 71,0 \%, no de melhor desempenho ( quarto superior ).
Nos quartos intermediários, essa proporção foi de 70,2 \% e 69,6 \% no segundo e no terceiro quartos de desempenho, respectivamente.
No Gráfico 4.6, pode -se constatar que não há diferença estatisticamente significativa entre as proporções de estudantes que consideraram a prova Longa ou Muito longa em função dos quartos de desempenho, sendo a maior proporção dessa resposta no primeiro quarto de desempenho ( 30,0 \% ), e a menor, no quarto superior ( 24,0 \% ), com uma tendência decrescente em relação ao desempenho.
Considerando -se a extensão da prova, o percentual de estudantes que avaliaram a prova como Longa ou Muito longa, em relação ao tempo total destinado à sua resolução foi maior para os estudantes de Centros Universitários, para os quais a proporção foi de 29,8 \%, enquanto as Universidades foram a Organização Acadêmica com a menor incidência, com 25,8 \%.
No Gráfico 4.11, é possível observar que não existem diferenças estatisticamente significativas entre as organizações.
Nas Organizações Acadêmicas, a proporção de presentes à prova que consideraram a sua extensão como Adequada, em relação ao tempo total destinado à sua resolução esteve entre 64,6 \%, nos Centros Universitários e 72,1 \% nas Universidades ( ver também Tabela II.6, no Anexo II ).
O percentual de estudantes que consideraram a prova como Longa ou Muito longa, em relação ao tempo total destinado à sua resolução foi maior para os estudantes de cursos de IES Privadas ( 27,4 \% ), sem diferença estatisticamente significativa da proporção em IES Públicas ( 26,3 \% ).
Já a proporção de presentes à prova que consideraram a extensão como Adequada foi 69,6 \% para os de IES Públicas e 68,1 \% para os de IES Privadas ( ver Gráfico 4.12 e Tabela II.6, no Anexo II ). 4.3 COMPREENSÃO DOS ENUNCIADOS DAS QUESTÕES 4.3.1 Componente de Formação Geral Em relação aos enunciados das questões do Componente de Formação Geral ( Questão 4 ), as opiniões foram positivas, já que 77,9 \% dos alunos avaliados consideraram todos ou a maioria dos enunciados das questões claros e objetivos ( Gráfico 4.13, Gráfico 4.14, e, no Anexo II, a Tabela II.7 ).
Na análise regional, a \textbf{percentagem} de estudantes que avaliaram que todos ou a maioria dos enunciados das questões do Componente de Formação Geral estavam claros e objetivos variou de 66,0 \%, na região Nordeste a 79,8 \%, na região Sudeste.
A análise das percepções dos estudantes sobre a clareza e a objetividade dos enunciados permite afirmar que todos, ou a maioria dos enunciados de questões relativas ao Componente de Formação Geral, foram considerados claros e objetivos para a maior parte dos respondentes ( maior ou igual a 66,0 \%, em todas as regiões e maior ou igual a 75,3 \% para todos os quartos de desempenho ).
Segundo o desempenho, observa -se que a proporção dos que emitiram essa opinião \textbf{cresce} num primeiro momento ; após isso, apresenta uma tendência decrescente, conforme o desempenho aumenta, sem diferenças estatisticamente significativas entre todos os quartos de desempenho.
No quarto superior, a clareza e a objetividade de todos ou da maioria dos enunciados das questões foram percebidas por 78,1 \%.
A proporção dos presentes que consideraram que todos ou a maioria dos enunciados das questões do Componente de Formação Geral estavam claros e objetivos foi maior para os estudantes das Faculdades, para os quais a proporção foi de 79,1 \%, enquanto os Centros Universitários foram a Organização Acadêmica com a menor incidência, com 74,7 \%.
No Gráfico 4.15, é possível observar que as diferenças entre todas as Organizações Acadêmicas não são estatisticamente significativas ( ver também Tabela II.8, no Anexo II ).
O percentual de estudantes que consideraram que todos ou a maioria dos enunciados das questões do Componente de Formação Geral estavam claros e objetivos foi maior para os estudantes de cursos de IES Privadas ( 79,5 \% ), sem diferença estatisticamente significativa da proporção em IES Públicas ( 75,3 \% ).
Um maior detalhamento pode ser encontrado no Gráfico 4.16 e na Tabela II.8, no Anexo II. 4.3.2 Componente de Conhecimento Específico Em relação aos enunciados das questões do Componente de Conhecimento Específico da prova, para 82,2 \% dos estudantes avaliados da Área de Tecnologia em Agronegócios, a clareza e a objetividade ( Questão 5 ) estavam presentes em todas ou na maioria das questões ( Gráfico 4.17 a Gráfico 4.20, e no Anexo II, as Tabelas II.9 e II.10 ).
A maioria dos estudantes de todas as Grandes Regiões brasileiras considerou claros e objetivos todos ou a maioria dos enunciados das questões do Componente de Conhecimento Específico da prova, percentual sempre maior ou igual a 74,0 \%.
Não existem diferenças significativas entre as regiões.
A proporção de estudantes que consideraram os enunciados das questões claros e objetivos apresenta uma tendência crescente em relação ao aumento de desempenho : mais elevada no quarto superior ( 88,3 \% ), se comparada ao quarto inferior de desempenho ( 77,3 \% ).
Não existem diferenças estatisticamente significativas entre os quartos de desempenho.
A proporção dos presentes que consideraram que todos ou a maioria dos enunciados das questões do Componente de Conhecimento Específico estavam claros e objetivos foi maior para os estudantes das Faculdades, para os quais a proporção foi de 84,9 \%, enquanto os Centros Universitários foram a Organização Acadêmica com a menor incidência, com 76,7 \%.
No Gráfico 4.19, é possível observar que não existem diferenças estatisticamente significativas entre as Organizações Acadêmicas ( ver também Tabela II.10, no Anexo II ).
O percentual de estudantes que consideraram que todos ou a maioria dos enunciados das questões do Componente de Conhecimento Específico estavam claros e objetivos foi maior para os estudantes de cursos de IES Públicas ( 83,4 \% ), sem diferença estatisticamente significativa da proporção em IES Privadas ( 80,5 \% ).
Ver Gráfico 4.20 e Tabela II.10, no Anexo II, para um maior detalhamento. 4.4 SUFICIÊNCIA DAS INFORMAÇÕES/INSTRUÇÕES FORNECIDAS Ao avaliarem as informações/instruções fornecidas para a resolução das questões ( Questão 6 ), 85,2 \% dos respondentes da Área de Tecnologia em Agronegócios de todo o Brasil afirmaram que estas eram até excessivas ou suficientes em todas ou na maioria das questões ( Gráfico 4.21, Gráfico 4.22, e, no Anexo II, a Tabela II.11 ).
Quanto à distribuição de respondentes pelas Grandes Regiões, observa -se que a proporção de estudantes que consideraram as informações/instruções fornecidas até excessivas ou suficientes em todas ou na maioria das questões foi sempre superior ou igual a 82,0 \%, chegando a 88,1 \%, na região Sul.
Não existem diferenças estatisticamente significativas entre as regiões.
Levando -se em conta o desempenho dos participantes, observa -se que não há um padrão para afirmar algum tipo de tendência das proporções de participantes que avaliaram as informações/instruções como até excessivas ou suficientes em todas ou na maioria das questões, havendo diferença estatisticamente significativa somente entre o quarto inferior ( 79,8 \% ) e o terceiro quarto ( 89,4 \% ) de desempenho, o percentual mais elevado, além de ser o percentual superior à média nacional ( 85,2 \% ), como mostrado no Gráfico 4.22.
Considerando -se as informações/instruções fornecidas para a resolução das questões ( Questão 6 ), a proporção que considerou que estas eram até excessivas ou suficientes em todas ou na maioria das questões foi maior para os estudantes das Faculdades, para os quais a proporção foi de 87,8 \%, enquanto as Universidades foram a Organização Acadêmica com a menor incidência, com 81,3 \%.
No Gráfico 4.23, é possível observar que somente há diferença estatisticamente significativa entre as Organizações Acadêmicas Universidades e Faculdades ( ver também Tabela II.12, no Anexo II ).
O percentual de estudantes que consideraram que as informações/instruções fornecidas para a resolução das questões eram até excessivas ou suficientes em todas ou na maioria das questões foi maior para os estudantes de cursos de IES Públicas ( 85,7 \% ), uma diferença estatisticamente significativa da proporção em IES Privadas ( 84,3 \% ), ver Gráfico 4.24 e Tabela II.12, no Anexo II. 4.5 DIFICULDADE ENCONTRADA AO RESPONDER À PROVA \textbf{Indagados} sobre as dificuldades com as quais se depararam ao responder à prova ( Questão 7 ), 15,0 \% dos estudantes apontaram o Desconhecimento do conteúdo.
Para 51,5 \%, a Forma diferente de abordagem do conteúdo foi indicada como dificuldade.
Já a Falta de motivação para fazer a prova foi a dificuldade apontada por 5,6 \% dos respondentes.
Considerando -se todo o Brasil, 25,0 \% dos respondentes afirmaram que não tiveram qualquer tipo de dificuldade para responder à prova ( Tabelas II.13 e II.14, no Anexo II ).
Nos Gráficos 4.25 a 4.28, são apresentados os percentuais de estudantes que apontaram o Desconhecimento do conteúdo como dificuldade ao responder à prova.
Na análise por Grandes Regiões, o percentual de inscritos e presentes que apontaram o Desconhecimento do conteúdo como dificuldade ao responder à prova não superou 17,4 \%.
Os percentuais variaram de 13,8 \%, na região Sudeste a 17,4 \%, na região Norte.
A Forma diferente de abordagem do conteúdo foi a escolha modal dos estudantes, com percentuais que variaram de 36,0 \% ( região Nordeste ) a 58,7 \% ( região Norte ).
O percentual de alunos que citaram a Falta de motivação como dificuldade variou de 4,0 \% ( região Nordeste ) a 8,7 \% ( região Norte ).
Os que declararam não ter qualquer tipo de dificuldade para responder à prova variaram de 13,0 \%, na região Norte a 38,0 \%, na região Nordeste.
Em relação aos quartos de desempenho, o Desconhecimento do conteúdo foi a opção escolhida por 16,9 \% dos estudantes do quarto inferior e por 16,8 \% do quarto superior e, apesar de uma variação para baixo, o gráfico não sugere um padrão para caracterizar uma tendência decrescente.
Não são observadas diferenças estatisticamente significativas entre todos os quartos de desempenho.
A alternativa modal para os alunos, quando agregados pelos quartos de desempenho, para a dificuldade encontrada, foi causada pela Forma diferente de abordagem do conteúdo : 51,4 \%, no quarto inferior e 54,3 \%, do terceiro quarto escolheram essa alternativa.
Na análise por Organização Acadêmica, o percentual de inscritos e presentes que apontaram o Desconhecimento do conteúdo como dificuldade ao responder à prova não superou 16,6 \%.
Os percentuais variaram de 13,7 \%, para CEFET/IF a 16,6 \%, nos Centros Universitários.
A Forma diferente de abordagem do conteúdo foi a escolha modal dos estudantes, com percentuais que variaram de 47,5 \% ( Universidades ) a 54,2 \% ( CEFET/IF ).
O percentual de alunos que citaram a Falta de motivação como dificuldade variou de 2,5 \% ( Centros Universitários ) a 8,8 \% ( CEFET/IF ).
Os que declararam não ter qualquer tipo de dificuldade para responder à prova variaram de 21,0 \%, nas CEFET/IF a 27,4 \%, nas Universidades.
No Gráfico 4.27, é possível observar que não existem diferenças estatisticamente significativas entre as Organizações Acadêmicas ( ver também Tabela II.14, no Anexo II ).
Em relação às Categorias Administrativas, o Desconhecimento do conteúdo foi a opção escolhida por 15,2 \% dos estudantes, nas IES Privadas e por 14,9 \% nas IES Públicas.
Não são observadas diferenças estatisticamente significativas entre estas duas categorias.
A alternativa modal para os alunos, quando agregados pelas Categorias Administrativas, para a dificuldade encontrada, foi causada pela Forma diferente de abordagem do conteúdo : 51,9 \%, para IES Públicas e 51,0 \%, para as Privadas ( ver Gráfico 4.28 e Tabela II.14, no Anexo II ). 4.6 CONTEÚDOS DAS QUESTÕES OBJETIVAS DA PROVA Ao analisar os conteúdos das questões objetivas da prova ( Questão 8 ), um percentual muito pequeno dos estudantes avaliados, apenas 3,4 \%, afirmaram que não estudou ainda a maioria desses conteúdos ( Gráficos 4.29 a Gráfico 4.32, e as Tabelas II.15 e II.16, no Anexo II ).
A maioria ( 75,6 \% ) afirmou ter estudado e aprendido muitos ou todos os conteúdos.
Na análise por Grande Região, a proporção de respondentes que escolheram a opção não estudou ainda a maioria desses conteúdos foi pequena.
Observa -se que, nas regiões Norte ( 12,8 \% ), Nordeste ( 4,0 \% ) e Sul ( 3,6 \% ), as proporções foram maiores do que a média nacional ( 3,4 \% ).
Nas regiões Sudeste e Centro-Oeste, as situações foram \textbf{inversas} : as proporções foram mais baixas ( 2,8 \% ) e ( 3,0 \% ), respectivamente.
Não são observadas diferenças estatisticamente significativas entre as regiões.
Em todas as regiões, a maioria dos presentes afirmou ter estudado e aprendido muitos ou todos os conteúdos, com proporções variando entre 44,7 \%, na região Norte a 78,6 \%, na região Sudeste.
Considerando -se, separadamente, as opiniões de estudantes dos quatro quartos de desempenho, observa -se que, no quarto inferior, 4,6 \% ofereceram como resposta que não estudou ainda a maioria desses conteúdos, sendo 2,2 \% os do quarto superior com a mesma resposta.
Não existem diferenças estatisticamente significativas entre os alunos que optaram por esse motivo de dificuldade e os quartos de desempenho.
Tendo -se em conta o quarto superior, 84,2 \% dos alunos afirmaram ter estudado e aprendido muitos ou todos os conteúdos.
No outro extremo, no primeiro quarto, 62,7 \% optaram pelas mesmas categorias.
As proporções são crescentes com o desempenho.
Na análise por Organização Acadêmica, a proporção de respondentes que escolheram a opção não estudou ainda a maioria desses conteúdos foi pequena.
Observa -se que, nos Centros Universitários ( 7,4 \% ) e CEFET/IF ( 4,9 \% ), as proporções foram maiores do que a média nacional ( 3,4 \% ).
Nas Universidades e Faculdades, as situações foram \textbf{inversas} : as proporções foram mais baixas ( 3,1 \% ) e ( 1,9 \% ), respectivamente.
No Gráfico 4.31, é possível observar que só há diferença estatisticamente significativa entre os Centros Universitários ( 7,4 \% ) e as Faculdades ( 1,9 \% ) ( ver também Tabela II.16, no Anexo II ).
Em todos os tipos de Organização Acadêmica, a maioria dos presentes afirmou ter estudado e aprendido muitos ou todos os conteúdos, com proporções variando entre 66,8 \%, nos CEFET/IF a 80,8 \%, nas Faculdades.
Gráfico 4.31 – Percentual de estudantes que informaram que não estudaram ainda a maioria desses conteúdos, por Organização Acadêmica – Enade/2019 – Tecnologia em Agronegócios Fonte : MEC/Inep/Daes - Enade/2019 Em relação às Categorias Administrativas, observa -se que, nas IES Privadas, 3,9 \% ofereceram como resposta que não estudou ainda a maioria desses conteúdos, sendo 3,1 \% nas IES Públicas com a mesma resposta.
A diferença entre os alunos que optaram por esse motivo de dificuldade nas duas categorias não é estatisticamente significativa.
Considerando -se as IES Privadas, 76,9 \% dos alunos afirmaram ter estudado e aprendido muitos ou todos os conteúdos.
Na outra categoria, 74,7 \% optaram pelas mesmas categorias ( ver Gráfico 4.32 e Tabela II.16, no Anexo II ). 4.7 TEMPO GASTO PARA CONCLUIR A PROVA Ao responderem sobre o tempo de conclusão da prova ( Questão 9 ), mais de que três quartos dos estudantes ( 76,2 \% ) afirmaram ter gasto Entre duas e quatro horas ( Gráfico 4.33 a Gráfico 4.36 e, no Anexo II, as Tabelas II.17 e II.18 ).
Considerando -se as cinco Grandes Regiões brasileiras, os percentuais dos que utilizaram Entre duas e quatro horas para finalizar a prova nas regiões Sudeste ( 75,6 \% ) e Sul ( 75,8 \% ) foram inferiores ao percentual nacional.
Nas demais Grandes Regiões, o percentual de alunos que \textbf{dispensaram} Entre duas e quatro horas para concluir a prova foi : 78,3 \%, na região Norte, 82,0 \%, na região Nordeste e 76,8 \%, na região Centro-Oeste, como mostrado no Gráfico 4.33.
Não existem diferenças estatisticamente significativas entre as regiões.
Analisadas as alternativas escolhidas pelos estudantes que se situam nos diferentes quartos de desempenho, observou -se uma tendência crescente : uma proporção maior de participantes no quarto superior declarou ter gasto Entre duas e quatro horas para concluir a prova quando comparada com a proporção de participantes nos quartos inferiores.
Existe diferença estatisticamente significativa entre quarto inferior e os demais quartos de desempenho.
Na análise por Organização Acadêmica, os percentuais dos que utilizaram Entre duas e quatro horas para finalizar a prova nas Universidades e nos CEFET/IF ( 74,2 \% em ambas ) foram inferiores ao percentual nacional.
Nos demais tipos de Organização Acadêmica, o percentual de alunos que \textbf{dispensaram} Entre duas e quatro horas para concluir a prova ficou 77,0 \%, nas Faculdades e 80,9 \%, nos Centros Universitários mostrado no Gráfico 4.35.
Não existem diferenças estatisticamente significativas entre as Organizações Acadêmicas.
Em relação às Categorias Administrativas, observa -se que, nas IES Privadas, 79,7 \% declararam ter gasto Entre duas e quatro horas para concluir a prova, sendo 73,8 \% nas IES Públicas com a mesma resposta.
A diferença entre os alunos que optaram por esse motivo de dificuldade nas duas categorias não é estatisticamente significativa ( ver Gráfico 4.36 e Tabela II.18, no Anexo II ).

% matched lemmas: arredondamento, crescer, disjunto, dispensar, estimador, estudantes, indagar, interseção, inverso, numérico, ordenar, parâmetro, percentagem, questionário, recorte, reprodução, subconjunto, vidar
\end{textsample}
